\input{../tex-inputs/latex-leseansicht-vorspann}

\section[Theodor von Sosnosky an Arthur Schnitzler, 14. 10. 1901]{L04321 Theodor von Sosnosky an Arthur Schnitzler, 14. 10. 1901}
\nopagebreak\mylabel{L04321v}
\rehead{ }\normalsize\beginnumbering\briefempfaengerindex{Schnitzler, Arthur@\textsc{Schnitzler, Arthur}!zzzSosnosky, Theodor von@\emph{von Theodor von Sosnosky}!1901-10-141@{14. 10. 1901}|(be}
\toendnotes[C]{\smallbreak\pagebreak[2]}
\correspDesc{Versand  durch Theodor von Sosnosky am 14. 10. 1901 in Alland
\newline{}Erhalt  durch Arthur Schnitzler im Zeitraum [15. 10. 1901 – 19. 10. 1901?] in Wien}\toendnotes[C]{\smallbreak}
\Standort{DLA, A:Schnitzler, HS.1985.1.4640.}
\physDesc{Brief, 2 Blätter, 6 Seiten, 4292 Zeichen
\newline{}Handschrift: blaue Tinte, deutsche Kurrent}\toendnotes[C]{\smallbreak}
\pstart
           \raggedleft{}{\pb}Alland\oindex{Alland@\textbf{Alland}, \emph{Verwaltungsgebiet}|pw}, 14./10 1901\pend
           
\pstart{}Sehr geehrter Herr Doktor!\pend\vspace{0.5em}
\pstart
           Von einem kurzen Wiener\oindex{Wien@\textbf{Wien}, \emph{Verwaltungsgebiet}|pw} Aufenthalte zurückgekehrt,
               fand ich \label{K_L04321-1v}\edtext{Ihren freundlichen Brief}{\lemma{\textnormal{\emph{Ihren freundlichen Brief}}}\Cendnote{\textnormal{XXXX Auszeichnungsfehler: Dokument L04235 nicht gefunden.}}}\label{K_L04321-1}, für den ich Ihnen
               verbindlichſt danke und der mich ſehr freut, weil ich daraus zu erſehen glaube, daß
               Sie \label{K_L04321-2v}\edtext{meinen Brief}{\lemma{\textnormal{\emph{meinen Brief}}}\Cendnote{\textnormal{XXXX Auszeichnungsfehler: Dokument L04233 nicht gefunden.}}}\label{K_L04321-2}{ }\introOben{}\strikeout{u. meinen Artikel\pwindex{Sosnosky, Theodor von 4.\,1.\,1866 Budapest – 3.\,2.\,1943 Wien@\textsc{Sosnosky, Theodor von} (4.\,1.\,1866 Budapest – 3.\,2.\,1943 Wien), \emph{Schriftsteller}!Fall Schnitzler. Eine unbefangene Betrachtung@\strich\emph{Der Fall Schnitzler. Eine unbefangene Betrachtung}|pwv}}\introOben{} ſo aufgeno{\geminationm}en haben, wie ich ihn gemeint
               habe.\pend
           
\pstart
           Daß der Ehrenrat\orgindex{[Ehrenrat der österreichisch-ungarischen Armee zur Veröffentlichung von Lieutenant Gustl, 1901]@[Ehrenrat der österreichisch-ungarischen Armee zur Veröffentlichung von Lieutenant Gustl, 1901]|pw} die Novelle\pwindex{Schnitzler, Arthur 15. 5. 1862 Wien – 21. 10. 1931 ebd.@\textsc{Schnitzler, Arthur} (15. 5. 1862 Wien – 21. 10. 1931 ebd.), \emph{Schriftsteller, Mediziner}!Lieutenant Gustl. Novelle@\strich\emph{Lieutenant Gustl. Novelle}|pwv} d. h. das in dieſer
                  vorhanden ſein \uline{ſollende} Vergehen gegen die
               Standesehre zum Subſtrat der Anklage, \textsc{resp}. des Urteils
               gemacht, war ſicher ein kraſſer \label{K_L04321-3v}\edtext{\begin{otherlanguage}{french}\textsc{faux pas}\end{otherlanguage}}{\lemma{\textnormal{\emph{faux pas}}}\Cendnote{\textnormal{französisch: Fehltritt, Fehler}}}\label{K_L04321-3};
               zum mindeſten hätte er dann die Stellen zitiren müſſen, aus denen dieſe Verletzung
               deutlich hervorgeht. Daß er dies nicht gethan hat, zeigt {\pb}daß
               er ſeiner Aufgabe, ein klares Urteil zu fällen, nicht gewachſen war. Er hat ſich
               dabei nur von \uline{Gefühlen} leiten laſſen, und die dürfen
               natürlich nicht genügen, über jemanden den Stab zu brechen. Daß dieſe Gefühle aber
               die richtigen waren, geht aus den \uline{von mir}{ }\label{K_L04321-4v}\edtext{zitirten Stellen}{\lemma{\textnormal{\emph{zitirten Stellen}}}\Cendnote{\textnormal{Sosnosky\pwindex{Sosnosky, Theodor von 4.\,1.\,1866 Budapest – 3.\,2.\,1943 Wien@\textsc{Sosnosky, Theodor von} (4.\,1.\,1866 Budapest – 3.\,2.\,1943 Wien), \emph{Schriftsteller}|pwk} zitiert zwei Passagen aus der Novelle\pwindex{Schnitzler, Arthur 15. 5. 1862 Wien – 21. 10. 1931 ebd.@\textsc{Schnitzler, Arthur} (15. 5. 1862 Wien – 21. 10. 1931 ebd.), \emph{Schriftsteller, Mediziner}!Lieutenant Gustl. Novelle@\strich\emph{Lieutenant Gustl. Novelle}|pwkv}, um die
                  Feindseligkeit gegenüber dem Militär zu belegen: »{\dots}\so{Die Frau von meinem Hauptmann}, das wär’ ja
                     doch keine anständige Frau {\dots} ich könnt’ schwören:
                     der Libitzky und der Wermutek und der schäbige Stellvertreter, der hat sie auch
                     gehabt {\dots} aber die \so{Frau Mannheimer}{ }{\dots} ja das wär’ was anders, das wär’ doch auch ein Umgang gewesen, das hätt’
                     einen beinah’ zu \so{einem andern Menschen gemacht} – da
                     hätt’ man doch noch \so{einen anderen Schliff} gekriegt
                        {\dots}« und »{\dots} am liebsten möchten sie (die Sozialisten) gleich's ganze Militär abschaffen;
                     aber \so{wer ihnen dann helfen möcht’, wenn die Chinesen über sie kommen, daran denken sie nicht}«, Hervorhebungen,
                  Auslassungen und Ergänzung nach Sosnoskys\pwindex{Sosnosky, Theodor von 4.\,1.\,1866 Budapest – 3.\,2.\,1943 Wien@\textsc{Sosnosky, Theodor von} (4.\,1.\,1866 Budapest – 3.\,2.\,1943 Wien), \emph{Schriftsteller}|pwk}{ }Text\pwindex{Sosnosky, Theodor von 4.\,1.\,1866 Budapest – 3.\,2.\,1943 Wien@\textsc{Sosnosky, Theodor von} (4.\,1.\,1866 Budapest – 3.\,2.\,1943 Wien), \emph{Schriftsteller}!Fall Schnitzler. Eine unbefangene Betrachtung@\strich\emph{Der Fall Schnitzler. Eine unbefangene Betrachtung}|pwkv}.}}}\label{K_L04321-4} wohl deutlich hervor. So ſchreibt nur ein
                  \uline{Feind} des Militärs! Aber ich will ſehr gerne
               zugeben, daß dieſe Stellen, ſelbst \uline{wenn}{ }ſie vom Ehrenrat\orgindex{[Ehrenrat der österreichisch-ungarischen Armee zur Veröffentlichung von Lieutenant Gustl, 1901]@[Ehrenrat der österreichisch-ungarischen Armee zur Veröffentlichung von Lieutenant Gustl, 1901]|pw} als \label{K_L04321-5v}\edtext{Gravamina}{\lemma{\textnormal{\emph{Gravamina}}}\Cendnote{\textnormal{lateinisch: Beschwerden}}}\label{K_L04321-5} angeführt
               worden wären, ein doch etwas zu dürftiges Subſtrat für eine Aberke{\geminationn}ung der \label{K_L04321-6v}\edtext{Charge}{\lemma{\textnormal{\emph{Charge}}}\Cendnote{\textnormal{obere Ranggruppe im
                  Militär}}}\label{K_L04321-6} gebildet hätten.\pend
           
\pstart
           Ich glaube übrigens, die ganze Sache hätte ein anderes Ende geno{\geminationm}en und Sie hätten ſich viel Unannehmlichkeiten erspart,
               wenn Sie der Aufforderung des Ehrenrats\orgindex{[Ehrenrat der österreichisch-ungarischen Armee zur Veröffentlichung von Lieutenant Gustl, 1901]@[Ehrenrat der österreichisch-ungarischen Armee zur Veröffentlichung von Lieutenant Gustl, 1901]|pw} sich zu
               rechtfertigen Folge {\pb}geleistet hätten, wozu Sie ja als Reſerve-Officier \uline{verpflichtet} waren. Zwiſchen \uline{Rechtfertigung} und \uline{Entschuldigung} ist
               ein großer Unterſchied! Hätten Sie vor dem Ehrenrat\orgindex{[Ehrenrat der österreichisch-ungarischen Armee zur Veröffentlichung von Lieutenant Gustl, 1901]@[Ehrenrat der österreichisch-ungarischen Armee zur Veröffentlichung von Lieutenant Gustl, 1901]|pw} erklärt, daß Ihnen die Ihnen \label{K_L04321-7v}\edtext{imputirte}{\lemma{\textnormal{\emph{imputirte}}}\Cendnote{\textnormal{veraltet: unterstellte, zur Last gelegte}}}\label{K_L04321-7} Absicht ferne gelegen ſei, daß
               Sie \strikeout{das} den Officiersſtand \uline{nicht} verhöhnen wollten,{ }ſo wäre das keine Entschuldigung, ſondern eine
               Rechtfertigung geweſen (ſofern das \uline{wirklich} nicht
               Ihre Absicht gewesen sein sollte), und ich glaube nicht, daß das Urteil dann ſo
               ausgefallen wäre, wie es thatſächlich ausgefallen ist.\pend
           
\pstart
           Von einer Selbsterniedrigung, einem Sich-Etwas-Vergeben konnte doch keine Rede{ }ſein; ſonst dürfte sich ja niemand, dem ein ſchwerer Vorwurf gemacht wird verteidigen
               und \substVorne{}\textsuperscript{man}\substDazwischen{}er\substHinten{} müßte auch dann ſchweigen, wenn das Civil{\pb}gericht, ihn zur
               Verantwortung zieht.\pend
           
\pstart
           Gegen gemeine Angriffe zu ſchweigen, das iſt ſicher die beſte und würdigſte Art der
               Abwehr; ſobald aber nicht private Feinde uns anklagen oder beſchuldigen, sondern
               irgend eine Behörde, der gegenüber wir die Verpflichtung haben, ihr \label{K_L04321-8v}\edtext{\textsc{Decorum}}{\lemma{\textnormal{\emph{Decorum}}}\Cendnote{\textnormal{lateinisch: das Schickliche}}}\label{K_L04321-8} zu
               wahren, und die daher das Recht hat, dies von uns zu verlangen, ſei es nun die
                  Ärzteka{\geminationm}er oder das betreffende Offizierscorps, ſo
               hieße das von Ihnen befolgte Prinzip nichts anderes als fataliſtiſche
               Selbstvernichtung. Die Logik iſt ja ſehr einfach; jede Behörde, ja jeder Einzelne muß
               so schließen: er hat Gelegenheit, ſich zu verteidigen, zu rechtfertigen; er thut es
               aber nicht, \uline{alſo} muß er ſchuldig ſein! Ein Mittelding
               gibts da nicht! Aus Ihrem Briefe geht ja auch ganz {\pb}deutlich hervor,
               daß Ihnen eine feindliche Abſicht \uline{nicht} ferne gelegen
               ist; denn Sie erklären, daß Sie – was Ihre Überzeugungstreue gewiß ſehr ehrt – unter
               allen Umſtänden die Wahrheit{ }ſagen wollen. D. h. Sie hatten die Frage des Ehrenrates\orgindex{[Ehrenrat der österreichisch-ungarischen Armee zur Veröffentlichung von Lieutenant Gustl, 1901]@[Ehrenrat der österreichisch-ungarischen Armee zur Veröffentlichung von Lieutenant Gustl, 1901]|pw}, ob Sie die Ihnen imputirte Absicht
               hatten, \uline{nicht} verneinen können, ohne die Unwahrheit
               zu ſagen.\pend
           
\pstart
           Ja, \uline{dann aber} ist das Urtheil nur gerecht, und
               Sie ſelber haben es dazu gemacht!\pend
           
\pstart
           Es fiele mir ſebſtverſtändlich nicht im Traume ein, Ihnen hier ſo unverhohlen meine
               Anſicht zu ſagen, die Sie ja nicht erbeten haben und daher als aufdringliche Anmaßung
               zurückweiſen können; aber da ich nun einmal öffentlich das Wort in dieſer Sache {\pb}ergriffen habe, ſo möchte ich auch privat nicht mißverstanden ſein,
               und ich glaube, daß mein Artikel\pwindex{Sosnosky, Theodor von 4.\,1.\,1866 Budapest – 3.\,2.\,1943 Wien@\textsc{Sosnosky, Theodor von} (4.\,1.\,1866 Budapest – 3.\,2.\,1943 Wien), \emph{Schriftsteller}!Fall Schnitzler. Eine unbefangene Betrachtung@\strich\emph{Der Fall Schnitzler. Eine unbefangene Betrachtung}|pwv} – vielleicht durch unklare Stiliſirung meinerſeits, von Ihrer Seite
               anders gedeutet worden iſt als er gemeint war. Daher dieſe Replik, für die ich ſonst
               keine Urſache hätte und ſomit auch keine Entſchuldigung verdiente, da ſie ohne
               Vorausgehen meines Artikels\pwindex{Sosnosky, Theodor von 4.\,1.\,1866 Budapest – 3.\,2.\,1943 Wien@\textsc{Sosnosky, Theodor von} (4.\,1.\,1866 Budapest – 3.\,2.\,1943 Wien), \emph{Schriftsteller}!Fall Schnitzler. Eine unbefangene Betrachtung@\strich\emph{Der Fall Schnitzler. Eine unbefangene Betrachtung}|pwv}
               eine plumpe, freche Taktloſigkeit wäre. Die aber liegt mir unendlich fern! Und eben
               dieſe Erwiderung scheint mir ein vielleicht ſonderbarer, aber ſicherer Beweis, daß
               mir an Ihrer Meinung etwas gelegen iſt und daß ich Ihrer Persönlichkeit ein
               beſonderes Intereße entgegenbringe. Indem ich hoffe, daß unsere nächste Correſpondenz
               wieder friedlicher und unverfänglicher werde, verbleibe ich, ſehr geehrter Herr
               Doktor, mit\pend
           \pstart vorzüglicher Hochachtung \spacefill\mbox{ThvSosnosky}\pend{}\selectlanguage{ngerman}\endnumbering\briefempfaengerindex{Schnitzler, Arthur@\textsc{Schnitzler, Arthur}!zzzSosnosky, Theodor von@\emph{von Theodor von Sosnosky}!1901-10-141@{14. 10. 1901}|)be}\mylabel{L04321h}
\begin{anhang}
\end{anhang}\newcommand{\dateiname}{L04321}\newcommand{\titel}{Theodor von Sosnosky an Arthur Schnitzler, 14. 10. 1901}\newcommand{\editorInnen}{Herausgegeben von Jahnke, SelmaMüller, Martin Anton}\input{../tex-inputs/latex-leseansicht-abspann}
